\usepackage[utf8]{inputenc}
\usepackage[T1]{fontenc}
\usepackage{lmodern}
\usepackage[leqno]{amsmath}
\usepackage{amssymb}
\usepackage[english,french]{babel}
\usepackage[a4paper,margin=30mm,marginparwidth=18mm]{geometry}
\usepackage{microtype}
\usepackage{xcolor}
\usepackage{paracol}
\usepackage{fancyhdr}
\usepackage[hypertexnames=false]{hyperref}
\usepackage{bookmark}

\definecolor{TohokuBlue}{HTML}{17365D}
\definecolor{TohokuGray}{HTML}{667085}
\hypersetup{
  colorlinks=true,
  linkcolor=TohokuBlue,
  urlcolor=TohokuBlue,
  citecolor=TohokuBlue,
  pdfauthor={Alexandre Grothendieck},
  unicode=true
}

\setlength{\parindent}{1.35em}
\setlength{\parskip}{0.35em}
\setlength{\emergencystretch}{2em}
\pagestyle{fancy}
\fancyhf{}
\fancyhead[LE,RO]{\small\textsc{Sur quelques points d'algebre homologique}}
\fancyhead[RE,LO]{\small\nouppercase{\leftmark}}
\fancyfoot[C]{\thepage}

\newcommand{\TohokuTarget}[1]{\hypertarget{#1}{}}
\providecommand{\LayerTarget}[1]{\TohokuTarget{#1}}
\providecommand{\LayerContents}[1]{\addcontentsline{toc}{chapter}{#1}}
\providecommand{\SectionHeading}[1]{\noindent\textbf{#1}}
\providecommand{\SourcePage}[1]{%
  \TohokuTarget{tohoku:source-page:#1}%
  \marginpar{\raggedleft\scriptsize\color{TohokuGray}source p.~#1}%
}
\newcommand{\TohokuContinuation}[1]{%
  \par\smallskip\noindent\hfill
  {\scriptsize\itshape\color{TohokuGray}Continued on source p.~#1.}%
}
\newcommand{\TohokuPartBoundary}[1]{}
\newcommand{\TohokuTerminal}[1]{}
\newcommand{\TohokuAuthorityLine}{%
  Tohoku Mathematical Journal (Second Series), vol.~9 (1957), pp.~119--221%
}
